% GENERATED FILE. Edit canonical lessons, then run npm run generate:books.
% canonical-source-hash: fnv1a64:36d038f3fc8e0d15
% canonical-lessons: IT-C404-bloccarsi, IT-C404-riavviare, IT-C404-virus, IT-C404-portatile, IT-C404-antivirus, IT-R404-ripasso-403-404, IT-R404-ripasso-final-404

\chapter{Non funziona}
\label{ch:it-non-funziona}
\hlchaptermodality{404}

\begin{chapteropening}
\textbf{By the end of this chapter:} \emph{I can say what is wrong with a device: it freezes, it needs restarting, it has a virus, and I need an antivirus.}

Complete the last lesson of chapter 404: I can say what is wrong with a device: it freezes, it needs restarting, it has a virus, and I need an antivirus.
\end{chapteropening}

\section[bloccarsi]{bloccarsi — to freeze, to get stuck}

\label{lesson:IT-C404-bloccarsi}



\begin{quote}
Before the new one: say the Italian for the email, then the Italian for the inbox.
\end{quote}

\subsection*{You'll want to know: bloccarsi}
\textbf{bloccarsi} — \textquotedblleft{}to freeze, to get stuck\textquotedblright{}.

\begin{quote}
Se il computer si blocca, aspetta un minuto.
\end{quote}

\begin{quote}
\emph{If the computer freezes, wait a minute.}
\end{quote}

Of a door or a car too: \emph{la porta si è bloccata}.

\begin{cousinweb}
\emph{bloccarsi} is built on \emph{blocco} \textquotedblleft{}block\textquotedblright{}, which was borrowed from French \emph{bloc} \textquotedblleft{}a heavy piece, a block\textquotedblright{} (EtymDB), from Middle Dutch \emph{blok} \textquotedblleft{}a tree trunk, a log\textquotedblright{}.
\end{cousinweb}

\begin{grammarlens}[title={What you've built}]
One more word for telling, explaining and describing at B1.
\end{grammarlens}

\subsection*{Your turn}
\begin{itemize}
\raggedright
  \item \emph{Say it:} \emph{bloccarsi}
  \item \emph{Say it:} \emph{bloccarsi}, once more
  \item \emph{Say it:} a sentence of your own with \emph{bloccarsi}
  \item \emph{Recall:} say the Italian for the email, then the Italian for the inbox, then say \emph{bloccarsi} again
\end{itemize}

\begin{tcolorbox}[breakable,colback=teal!4,colframe=teal!35!black,title={Before you move on}]
What does bloccarsi mean? (\textquotedblleft{}To freeze, to get stuck\textquotedblright{}.) Say it once more.
\end{tcolorbox}
\section[riavviare]{riavviare — to restart, to reboot}

\label{lesson:IT-C404-riavviare}



\begin{quote}
Before the new one: say the Italian for the inbox, then the Italian for to freeze.
\end{quote}

\subsection*{You'll want to know: riavviare}
\textbf{riavviare} — \textquotedblleft{}to restart, to reboot\textquotedblright{}.

\begin{quote}
Se non funziona, riavvia il telefono.
\end{quote}

\begin{quote}
\emph{If it does not work, restart the phone.}
\end{quote}

\emph{Avviare} is to start; \emph{ri-} is again.

\begin{cousinweb}
\emph{riavviare} is built from \emph{ri-} \textquotedblleft{}again\textquotedblright{} and \emph{avviare} \textquotedblleft{}to start off\textquotedblright{}, from \emph{a} and \emph{via} \textquotedblleft{}way, road\textquotedblright{}, inherited from Latin \emph{via} (Lewis \& Short): to set something on its way again.
\end{cousinweb}

\begin{grammarlens}[title={What you've built}]
One more word for telling, explaining and describing at B1.
\end{grammarlens}

\subsection*{Your turn}
\begin{itemize}
\raggedright
  \item \emph{Say it:} \emph{riavviare}
  \item \emph{Say it:} \emph{riavviare}, once more
  \item \emph{Say it:} a sentence of your own with \emph{riavviare}
  \item \emph{Recall:} say the Italian for the inbox, then the Italian for to freeze, then say \emph{riavviare} again
\end{itemize}

\begin{tcolorbox}[breakable,colback=teal!4,colframe=teal!35!black,title={Before you move on}]
What does riavviare mean? (\textquotedblleft{}To restart, to reboot\textquotedblright{}.) Say it once more.
\end{tcolorbox}
\section[il virus]{il virus — the virus}

\label{lesson:IT-C404-virus}



\begin{quote}
Before the new one: say the Italian for to freeze, then the Italian for to restart.
\end{quote}

\subsection*{You'll want to know: il virus}
\textbf{il virus} — \textquotedblleft{}the virus\textquotedblright{}.

\begin{quote}
Se apri quell'allegato, prendi un virus.
\end{quote}

\begin{quote}
\emph{If you open that attachment, you will get a virus.}
\end{quote}

For computers and for illness. Masculine, invariable.

\begin{cousinweb}
\emph{virus} is taken whole from Latin \emph{vīrus} \textquotedblleft{}slime, poison\textquotedblright{} (Lewis \& Short); the computer sense came from English.
\end{cousinweb}

\begin{grammarlens}[title={What you've built}]
One more word for telling, explaining and describing at B1.
\end{grammarlens}

\subsection*{Your turn}
\begin{itemize}
\raggedright
  \item \emph{Say it:} \emph{il virus}
  \item \emph{Say it:} \emph{il virus}, once more
  \item \emph{Say it:} a sentence of your own with \emph{il virus}
  \item \emph{Recall:} say the Italian for to freeze, then the Italian for to restart, then say \emph{il virus} again
\end{itemize}

\begin{tcolorbox}[breakable,colback=teal!4,colframe=teal!35!black,title={Before you move on}]
What does il virus mean? (\textquotedblleft{}The virus\textquotedblright{}.) Say it once more.
\end{tcolorbox}
\section[il portatile]{il portatile — the laptop}

\label{lesson:IT-C404-portatile}



\begin{quote}
Before the new one: say the Italian for to restart, then the Italian for the virus.
\end{quote}

\subsection*{You'll want to know: il portatile}
\textbf{il portatile} — \textquotedblleft{}the laptop\textquotedblright{}.

\begin{quote}
Porto il portatile in vacanza, se devo lavorare.
\end{quote}

\begin{quote}
\emph{I take my laptop on holiday, if I have to work.}
\end{quote}

As an adjective, portable: \emph{un televisore portatile}.

\begin{cousinweb}
\emph{portatile} is built on \emph{portare} \textquotedblleft{}to carry\textquotedblright{}, inherited from Latin \emph{portāre} \textquotedblleft{}to carry\textquotedblright{} (Lewis \& Short): a computer you can carry.
\end{cousinweb}

\begin{grammarlens}[title={What you've built}]
One more word for telling, explaining and describing at B1.
\end{grammarlens}

\subsection*{Your turn}
\begin{itemize}
\raggedright
  \item \emph{Say it:} \emph{il portatile}
  \item \emph{Say it:} \emph{il portatile}, once more
  \item \emph{Say it:} a sentence of your own with \emph{il portatile}
  \item \emph{Recall:} say the Italian for to restart, then the Italian for the virus, then say \emph{il portatile} again
\end{itemize}

\begin{tcolorbox}[breakable,colback=teal!4,colframe=teal!35!black,title={Before you move on}]
What does il portatile mean? (\textquotedblleft{}The laptop\textquotedblright{}.) Say it once more.
\end{tcolorbox}
\section[l'antivirus]{l'antivirus — the antivirus (software)}

\label{lesson:IT-C404-antivirus}



\begin{quote}
Before the new one: say the Italian for the virus, then the Italian for the laptop.
\end{quote}

\subsection*{You'll want to know: l'antivirus}
\textbf{l'antivirus} — \textquotedblleft{}the antivirus (software)\textquotedblright{}.

\begin{quote}
Se hai un buon antivirus, sei più tranquillo.
\end{quote}

\begin{quote}
\emph{If you have a good antivirus, you can feel calmer.}
\end{quote}

Masculine, invariable: \emph{l'antivirus, gli antivirus}.

\begin{cousinweb}
\emph{antivirus} is built from the Greek prefix \emph{anti-} \textquotedblleft{}against\textquotedblright{} and \emph{virus}, from Latin \emph{vīrus} \textquotedblleft{}slime, poison\textquotedblright{} (Lewis \& Short).
\end{cousinweb}

\begin{grammarlens}[title={What you've built}]
One more word for telling, explaining and describing at B1.
\end{grammarlens}

\subsection*{Your turn}
\begin{itemize}
\raggedright
  \item \emph{Say it:} \emph{l'antivirus}
  \item \emph{Say it:} \emph{l'antivirus}, once more
  \item \emph{Say it:} a sentence of your own with \emph{l'antivirus}
  \item \emph{Recall:} say the Italian for the virus, then the Italian for the laptop, then say \emph{l'antivirus} again
\end{itemize}

\begin{tcolorbox}[breakable,colback=teal!4,colframe=teal!35!black,title={Before you move on}]
What does l'antivirus mean? (\textquotedblleft{}The antivirus (software)\textquotedblright{}.) Say it once more.
\end{tcolorbox}
\section[(dialogue)]{Condizioni e regole}

\label{lesson:IT-R404-ripasso-403-404}



\begin{quote}
Say the Italian for to publish and to comment.
\end{quote}

\subsection*{Your turn}
Say these aloud:

\begin{itemize}
\raggedright
  \item to publish, to comment, to forward, the email, the inbox
  \item to freeze, to restart, the virus, the laptop, the antivirus
\end{itemize}

\begin{tcolorbox}[breakable,colback=teal!4,colframe=teal!35!black,title={Before you move on}]
To publish? (\textbf{pubblicare}.) To comment? (\textbf{commentare}.) To forward? (\textbf{inoltrare}.) The email? (\textbf{l'email}.) The inbox? (\textbf{la casella}.) To freeze? (\textbf{bloccarsi}.) To restart? (\textbf{riavviare}.) The virus? (\textbf{il virus}.) The laptop? (\textbf{il portatile}.) The antivirus? (\textbf{l'antivirus}.)
\end{tcolorbox}
\section[(dialogue)]{Non funziona — once more}

\label{lesson:IT-R404-ripasso-final-404}



\begin{quote}
Say the Italian for to freeze and to restart.
\end{quote}

\subsection*{Your turn}
\begin{itemize}
\raggedright
  \item \emph{Say it:} to freeze, to restart, the virus, the laptop, the antivirus
\end{itemize}

\begin{tcolorbox}[breakable,colback=teal!4,colframe=teal!35!black,title={Before you move on}]
To freeze? (\textbf{bloccarsi}.) To restart? (\textbf{riavviare}.) The virus? (\textbf{il virus}.) The laptop? (\textbf{il portatile}.) The antivirus? (\textbf{l'antivirus}.)
\end{tcolorbox}
